% !TeX root = ../main.tex
\section*{Summary}
\addcontentsline{toc}{section}{Summary}

Preference-Tree Optimization (PTO) trains a small therapist language model by branching at
each therapist turn, simulating $\K$ further turns of the conversation under the current
policy, and letting an LLM oracle score the resulting trajectory rather than the reply in
isolation. Our ICLR 2025 workshop paper reported that this look-ahead was the lever that
mattered: deeper look-ahead produced higher and more stable scores
\citep{baruch2025pto}. This chapter asks whether that finding holds, using three
generations of the same experiment --- a published one and two subsequent
re-implementations --- and re-analysing all three with one set of statistics and one
measurement axis.

It does not hold. Persona-paired at matched iteration counts and graded throughout by the
same \primaryjudge{} oracle, look-ahead ($\K{=}5$ versus $\K{=}0$) moves from a clear
advantage in generation~1 (ahead at $7/7$ matched iterations, mean $\dlt = -0.132$ on
\QQ{}), through a clean null in generation~2 ($0$ of $105$ paired contrasts significant;
$\K{=}5$ ahead in $58.1\%$ of them, a coin flip), to never leading in generation~3. There
the held-out grader puts look-ahead behind at all $8$ matched iterations, and the primary
oracle at $7$ of $8$ --- the eighth being a $-0.002$ dead heat rather than a win --- while
one iteration significantly \emph{favours} $\K{=}0$ under both.

Two controls make this more than a list of disagreeing experiments. First, the reversal is
not an artifact of the statistics: re-analysing the published experiment with the modern
paired test makes its look-ahead effect \emph{stronger}, not weaker, than the unpaired
best-versus-best comparison originally reported. Second, it is not an artifact of the
grader: re-scoring all $1{,}344$ of generation~1's conversations with the exact oracle used
in generation~3 preserves the effect ($7/7$, $20$ of $21$ metric cells), so the judge
change cannot explain why the effect later disappears.

What does distinguish the generations is how well the \emph{myopic} baseline learns. Under
the common grader, generation~1's $\K{=}0$ arm never departs from the untrained base
($+0.023$ over seven iterations), whereas generation~3's climbs $+1.212$. Look-ahead's
measured benefit runs opposite to that learning rate across the five training arms for
which both quantities exist ($r = -0.72$ on the slope, $r = -0.84$ on the total gain) ---
suggestive, but with $n = 5$ arms and one arm contradicting it, an interpretation rather
than a result. Three independent measurements on the training signal itself agree with the
reading: look-ahead multiplies the spread of candidate rewards by ${\sim}1.55\times$ with
margin and standard deviation rising by the \emph{identical} factor, creates no additional
branch points, and at matched policy adds no measurable reward faithfulness.

The practical conclusion is that look-ahead simulation is a rescue for a reward signal that
is failing, not a general-purpose improvement --- and that in this domain it costs several
times more per iteration to buy nothing once the one-step signal works.
